\documentclass[10pt,a4paper]{amsart}
\usepackage[margin=19mm]{geometry}
\usepackage{amsmath,amssymb,mathtools}
\usepackage[scr=rsfs,cal=euler]{mathalfa}
\usepackage{xfrac}
\usepackage[unicode,hidelinks,bookmarks=false]{hyperref}
\newcommand{\ie}{i.e.\ }
\newcommand{\cf}{cf.\ }
\newcommand{\resp}{respectively\ }
\newcommand{\Subsection}{\subsection}
\providecommand{\nobreakdash}{\nobreak\hbox{-}}
\setlength{\emergencystretch}{2em}
\multlinegap0pt
\theoremstyle{plain}
\newtheorem{Thm}{Theorem}
\newtheorem{Cor}[Thm]{Corollary}
\newtheorem{Claim}[Thm]{Claim}
\newtheorem{lemma}[Thm]{Lemma}
\newtheorem{Prop}[Thm]{Proposition}
\newtheorem{Rem}[Thm]{Remark}
\title[Stability of fixed life histories to per]{Stability of fixed life histories to perturbation by rare diapause: a packing bound for the family of strings of length $T$ with $k$ changes}
\author{David Steinsaltz; Shripad Tuljapurkar}
\date{Source: arXiv:1809.04027v3 (CC BY 4.0)}
\usepackage{euscript}
\ifdefined\F\else
\newcommand{\F}{\eu{F}}
\fi
\ifdefined\e\else
\newcommand{\e}{\mathrm{e}}
\fi
\ifdefined\eu\else
\newcommand{\eu}{\EuScript}
\fi
\begin{document}
\maketitle
\noindent\emph{Source and licence.} Stability of fixed life histories to perturbation by rare diapause. Version arXiv:1809.04027v3 (CC BY 4.0), distributed under the Creative Commons Attribution 4.0 International licence, \url{http://creativecommons.org/licenses/by/4.0/}, as recorded in the arXiv licence metadata for this exact version and shown on the abstract page \url{https://arxiv.org/abs/1809.04027v3}. Full licence notice: licenses/arxiv-1809.04027-CC-BY-4.0.txt.\par\smallskip

\noindent\textit{Editorial scope.} Counted unit: the article's Lemma with source label \texttt{L:packingbound} (source0.tex lines 1211--1220), the packing bound for the family of strings of length $T$ over an alphabet of size $d$ with $k$ changes, given both for radii above the critical value and for all radii; delivered together with the article's complete original proof of it (source0.tex lines 1222--1252, one \texttt{proof} environment of 31 lines, adjacent to the statement). No mathematics is written, repaired or re-derived: statement and proof are the article's own text line for line, in the article's own notation (the family of strings, the change set, the packing number, the entropy-type bounds), and no retained line is substituted, re-flowed or omitted. No further part of the article is required as context and none is retained: the counted proof uses only the objects introduced in its own statement together with the elementary binomial estimate that it states, so no cross-reference and no citation occurs in any retained line; consequently the wrapper carries no bibliography and no bibliography entry, and no context passage is delivered. Nothing else of the article is retained: its main theorems, its other lemmas and its figures are omitted. Editorial formatting only: an AMS \texttt{amsart} preamble that restates the environments and the macros the retained lines use, namely the article's declarations of Theorem, Corollary, Claim, Lemma, Proposition and Remark sharing one master counter, and the macros for the special letters and operators of the retained lines. The article's own source declares the class \texttt{amsart} as well; no class or style file is shipped with this package. Numbering differs from the article, because the article declares its Lemma environment as sharing the Theorem master counter and only a subset is retained: the counted lemma prints with the wrapper's own number. No picture environment and no graphical inclusion occurs in any retained line, and the package delivers no image asset. One bundled source file, \texttt{sources/101814/source0.tex}, accompanies the delivered item as the exact source of every retained span. Every retained span is recorded in \texttt{delivery-evidence.json} with method \texttt{exact\_tex}.
\begin{lemma} \label{L:packingbound}
For positive integers $T,k$, and any positive $r$ with $\tau\sqrt{T} \ge r > \tau \sqrt{k}$,
\begin{equation}  \label{E:packingbound}
D(r \, ;\, \F_{T,k})\le d^{k+1} \left( \frac{T\e}{(r/\tau)^{2}-k}+2\e \right)^{k}.
\end{equation}
and for all $r$
\begin{equation} \label{E:packingbound2}
D(r \, ;\, \F_{T,k}) \le d^{k+1}\left( \frac{T\e}{k} \right)^k
\end{equation}
\end{lemma}

\begin{proof}
Let $r'=(r/\tau)^{2}$. Suppose that $r'>k$, let $j=\lfloor r'/k\rfloor$, and let $m=\lceil T/j\rceil$. Let $\F_{*}$ be the set of ordered (non-decreasing) sequences
of length $k$ from $\{0,\dots,m-1\}$, crossed with $\{0,\dots,d-1\}^{k+1}$, 
and define a map $(\phi,\psi):\F_{T,k}\to \F_{*}$ 
by letting $\{f\}_{i}$ be the $i$-th coordinate where $f$ changes, and defining
$$
\phi(f)_{i}= \bigl\lfloor \frac{\{f\}_{i}}{j}  \bigr\rfloor
$$
and $\psi(f)_{i}=f_{\{f\}_{i}}$; that is, the site that $f$ moves to at its $i$-th
change.

If $f$ and $f'$ are two elements of $\F_{T,k}$ with $\phi(f)=\phi(f')$
and $\psi(f)=\psi(f')$, 
then $f_{i}=f'_{i}$ as long as $\lfloor i/j\rfloor\notin \phi(f)$, 
since any $t\notin \phi(f)$ corresponds to a span of $tj,tj+1,\dots,tj+j-1$
where $f_{tj}=f'_{tj}$ (because they started with $f_{0}=f'_{0}$, and the number
of changes in $f_{0},\dots,f_{tj-1}$ is the same as the number
of changes in $f'_{0},\dots,f'_{ti-1}$). 
Thus the Hamming distance is no more than $k\cdot j \le r'$, meaning that $\|f-f'\|_{2}\le\sqrt{r'}$, so $\rho(f,f')\le \tau\sqrt{r'}=r$. By the
pigeonhole principle, any subset of $\F_{T,k}$ of size greater than $\# \F_{*}$
has points with $\rho$-separation no more than $r$. Hence
$$
D(r,\F_{T,k})\le \#\F_{*} \le \binom{m+k}{k} d^{k+1}.
$$

Combining this with the trivial bound $D(r,\F_{T,k})\le \#\F_{T,k}\le d^{k+1}\binom{T}{k}$ and the bound
$$
\binom{a}{b}\le \left(\frac{a\e}{b}\right)^{b}
$$
completes the proof.
\end{proof}

\end{document}
